% !TEX root = ../Main.tex
\chapter{传送带}

传送带问题的关键只有两点：摩擦力的方向由\textbf{物块相对传送带}的运动方向决定（与之相反），而不是相对地面；以及物块和传送带\textbf{共速}的那一刻要做检验——共速之后摩擦力是否还能维持共速、会不会改向。下面分水平、倾斜两类，逐一画 $v$-$t$ 图像。

\section{水平传送带}

把物块由静止放到一条向右匀速（速度 $v_0$）的传送带上。

\begin{center}
\begin{tikzpicture}[>=Stealth, scale=1]
    \draw[thick] (-2.3,0) -- (2.3,0);
    \draw[thick] (-2.6,0) circle (0.28);
    \draw[thick] (2.6,0) circle (0.28);
    \draw[thick, fill=gray!15] (-0.45,0) rectangle (0.45,0.7);
    \draw[->, thick] (0.7,0.35) -- (1.7,0.35) node[right] {\small $v_0$};
\end{tikzpicture}
\end{center}

物块初速为零、传送带向右，物块相对传送带向左，所以摩擦力向右，物块以 $a=\mu g$ 加速。

\rmkb[注意]{
    \textbf{共速检验。} 物块加速到 $v_0$ 时与传送带共速。此后水平方向没有别的力，要维持共速只需静摩擦为零——能做到，于是物块与传送带一起匀速向右。$v$-$t$ 图像就是先一段斜直线（斜率 $\mu g$），到 $v_0$ 后转成水平线。
}

\section{倾斜传送带（传送带足够长）}

倾斜时多了重力沿斜面的分量，要和摩擦力一起算，但摩擦力方向的判据不变——还是看物块相对传送带往哪动。下面是两种典型情形，运动都分成两个阶段：

\begin{center}
\begin{tikzpicture}[>=Stealth, scale=0.92]
    % 情形一：传送带向下，物块从上端放下
    \begin{scope}
        \draw[thick] (0,0) -- (4.2,2.1);
        \draw[thick] (-0.25,-0.12) circle (0.2);
        \draw[thick] (4.45,2.22) circle (0.2);
        \draw[thick, fill=gray!15] (2.9,1.45) ++(27:-0.0) -- ++(27:0.55) -- ++(117:0.45) -- ++(27:-0.55) -- cycle;
        \draw[->, thick] (1.6,0.8) -- (0.7,0.35) node[below right] {\small $v_0$};
        \draw (0.55,0) arc (0:27:0.55);
        \node at (0.75,0.16) {\scriptsize $\theta$};
        \node at (2.1,-0.55) {\small 情形一：带向下、物块放在上端};
    \end{scope}
    % 情形二：传送带向上，物块从下端放下
    \begin{scope}[xshift=7cm]
        \draw[thick] (0,0) -- (4.2,2.1);
        \draw[thick] (-0.25,-0.12) circle (0.2);
        \draw[thick] (4.45,2.22) circle (0.2);
        \draw[thick, fill=gray!15] (1.0,0.5) -- ++(27:0.55) -- ++(117:0.45) -- ++(27:-0.55) -- cycle;
        \draw[->, thick] (2.4,1.2) -- (3.3,1.65) node[above left] {\small $v_0$};
        \draw (0.55,0) arc (0:27:0.55);
        \node at (0.75,0.16) {\scriptsize $\theta$};
        \node at (2.1,-0.55) {\small 情形二：带向上、物块放在下端};
    \end{scope}
\end{tikzpicture}
\end{center}

\state{两个阶段}{
    \textbf{第一阶段}：物块相对传送带滑动，摩擦力沿传送带运动方向（即把物块往带的方向拖），先把物块加速到与传送带共速。

    \textbf{第二阶段}（共速检验）：共速后看 $\mu$ 与 $\tan\theta$ 的关系——
    \begin{itemize}
        \item 若 $\mu\ge\tan\theta$，静摩擦足以抵住重力分量，物块与传送带一起匀速；
        \item 若 $\mu<\tan\theta$，静摩擦撑不住，物块继续相对带下滑，摩擦力改向沿斜面向上，物块以 $a=g(\sin\theta-\mu\cos\theta)$ 继续加速。
    \end{itemize}
}

两种情形的运动都分两段，下面分别画出 $v$-$t$ 图像（情形一取 $\mu<\tan\theta$、共速后继续下滑，情形二取 $\mu\ge\tan\theta$、共速后随带匀速）。

\begin{center}
\begin{tikzpicture}[>=Stealth, scale=1]
    % 情形一：带向下、物块放上端
    \begin{scope}
        \draw[->] (-0.3,0) -- (4.2,0) node[right] {\small $t$};
        \draw[->] (0,-0.3) -- (0,3.0) node[above] {\small $v$};
        % 第一阶段：a1=g(sinθ+μcosθ) 较陡，到共速 v0
        \draw[thick, blue] (0,0) -- (1.3,1.6);
        % 第二阶段：a2=g(sinθ-μcosθ) 较缓，继续加速
        \draw[thick, blue] (1.3,1.6) -- (3.6,2.6);
        \draw[dashed] (1.3,0) -- (1.3,1.6) -- (0,1.6);
        \node[left] at (0,1.6) {\small $v_0$};
        \node[below] at (1.3,-0.05) {\small $t_1$};
        \node at (2.1,-0.7) {\small 情形一：带向下、物块放上端};
    \end{scope}
    % 情形二：带向上、物块放下端
    \begin{scope}[xshift=6.5cm]
        \draw[->] (-0.3,0) -- (4.2,0) node[right] {\small $t$};
        \draw[->] (0,-0.3) -- (0,3.0) node[above] {\small $v$};
        % 第一阶段：a=g(μcosθ-sinθ) 加速到共速 v0
        \draw[thick, blue] (0,0) -- (1.6,1.6);
        % 第二阶段：共速后跟带匀速（水平）
        \draw[thick, blue] (1.6,1.6) -- (3.6,1.6);
        \draw[dashed] (1.6,0) -- (1.6,1.6) -- (0,1.6);
        \node[left] at (0,1.6) {\small $v_0$};
        \node[below] at (1.6,-0.05) {\small $t_1$};
        \node at (2.1,-0.7) {\small 情形二：带向上、物块放下端};
    \end{scope}
\end{tikzpicture}
\end{center}

\section{物块冲入传送带（取向左为正）}

水平传送带向右以 $v_2$ 运动，物块以 $v_1$ 从右端向左冲入，与带的动摩擦因数为 $\mu$。规定向左为正。

\begin{center}
\begin{tikzpicture}[>=Stealth, scale=1]
    \draw[thick] (-2.6,0) -- (2.6,0);
    \draw[thick] (-2.9,0) circle (0.26);
    \draw[thick] (2.9,0) circle (0.26);
    \draw[thick, fill=gray!15] (0.6,0) rectangle (1.4,0.65);
    \draw[->, thick] (0.5,0.32) -- (-0.5,0.32) node[left] {\small $v_1$};
    \draw[->, thick] (-1.3,-0.4) -- (-0.2,-0.4) node[right] {\small $v_2$};
\end{tikzpicture}
\end{center}

物块向左（正向）、传送带向右（负向），物块相对传送带向正向，摩擦力沿负向（向右），物块向左做减速。减速到零后被摩擦力反向加速（向右、即负向），直到与传送带共速 $-v_2$，之后一起向右匀速。$v$-$t$ 图像都是一条斜率为 $-\mu g$ 的斜线，过零后继续下降到 $-v_2$、再转为水平；按初速 $v_1$ 与 $v_2$ 的大小关系，分两种情况各画一张：

\begin{center}
\begin{tikzpicture}[>=Stealth, scale=1]
    % (i) v1 <= v2：起点低于 -v2 那条线的高度
    \begin{scope}
        \draw[->] (0,-2.1) -- (0,1.6) node[above] {\small $v$};
        \draw[->] (-0.3,0) -- (4.2,0) node[right] {\small $t$};
        \draw[thick, blue] (0,0.8) -- (1.95,-1.7);
        \draw[thick, blue] (1.95,-1.7) -- (3.8,-1.7);
        \draw[dashed] (1.95,-1.7) -- (0,-1.7);
        \node[blue, left] at (0,0.8) {\small $v_1$};
        \node[blue, left] at (0,-1.7) {\small $-v_2$};
        \node at (1.9,-2.55) {\small (i)\,$v_1\le v_2$};
    \end{scope}
    % (ii) v1 > v2：起点高于 -v2 那条线的高度
    \begin{scope}[xshift=6cm]
        \draw[->] (0,-1.6) -- (0,2.1) node[above] {\small $v$};
        \draw[->] (-0.3,0) -- (4.2,0) node[right] {\small $t$};
        \draw[thick, blue] (0,1.7) -- (2.7,-1.0);
        \draw[thick, blue] (2.7,-1.0) -- (3.8,-1.0);
        \draw[dashed] (2.7,-1.0) -- (0,-1.0);
        \node[blue, left] at (0,1.7) {\small $v_1$};
        \node[blue, left] at (0,-1.0) {\small $-v_2$};
        \node at (1.9,-2.0) {\small (ii)\,$v_1> v_2$};
    \end{scope}
\end{tikzpicture}
\end{center}

\section{绳连物块、传送带提速}

\ex{传送带上的物块用绳连在左侧的墙上，绳的张力已达最大值（即物块已处于将滑未滑的临界，绳拉力与动摩擦力平衡）。现增大传送带的速度，求 $1\,\mathrm{s}$ 后物块的速度大小。}

\sol{
    增大传送带速度之前，绳的张力 $T$ 与物块所受摩擦力平衡，物块静止。增大传送带速度后，物块相对传送带的运动方向\textbf{没有改变}（物块仍相对带向后），所以摩擦力方向、大小都不变，仍是滑动摩擦力 $\mu mg$；绳的张力仍与之平衡。物块所受合力为零，依旧静止。

    所以 $1\,\mathrm{s}$ 后物块速度仍为 $0$。常见的错误是以为"带变快物块就跟着变快"——但有绳拴着，物块不可能被带走。
}
